\subsection{\texorpdfstring{PROPERTIES OF THE FUNCTION \(e^{z}\)}{PROPERTIES OF THE FUNCTION E TO THE Z}}

The fact that \(z \mapsto e^z\) is a homomorphism of \(\mathbf{C}\) onto \(\mathbf{C}^*\) may be expressed by the identities

\[
e ^ {z + z ^ {\prime}} = e ^ {z} e ^ {z ^ {\prime}}, \qquad e ^ {0} = 1, \qquad e ^ {- z} = 1 / e ^ {z}.\tag{22}
\]

By definition, one has, for every \(z = x + iy\) ,

\[
e ^ {x + i y} = e ^ {x} (\cos y + i \sin y)\tag{23}
\]

and since \(e^x > 0\) one sees that \(e^z\) has absolute value \(e^x\) and amplitude \(y\) (modulo \(2\pi\) ).

In particular, def. 2 (III, p.~98) gives

\[
\mathbf {e} (x) = e ^ {2 \pi i x}\tag{24}
\]

which permits us to write the formulae defining \(\cos x\) and \(\sin x\) in the form

\[
\cos x = \frac {1}{2} \left(e ^ {i x} + e ^ {- i x}\right), \quad \sin x = \frac {1}{2 i} \left(e ^ {i x} - e ^ {- i x}\right)\tag{25}
\]

(Euler's formulae).

Since \(2\pi\) is the principal period of \(\mathbf{e}(y/2\pi)\) , \(2\pi i\) is the principal period of \(e^{z}\) ; in other words, the group of periods of \(e^{z}\) is the set of numbers \(2n\pi i\) , where n runs through Z.

Finally, formula (21) of III, p.~98 can be written

\[
\mathrm{D} \left(e ^ {z}\right) = e ^ {z}\tag{26}
\]

whence, for every complex number \(a\)

\[
\mathrm{D} \big (e ^ {a z} \big) = a e ^ {a z}.\tag{27}
\]

Remark. If, in formula (27), one restricts the function \(e^{az}\) (a complex) to the real axis, one again obtains, for x real,

\[
\mathrm{D} (e ^ {a x}) = a e ^ {a x}.\tag{28}
\]

This formula allows us to calculate a primitive for each of the functions \(e^{\alpha x} \cos \beta x\) , \(e^{\alpha x} \sin \beta x\) ( \(\alpha\) and \(\beta\) real); indeed we have \(e^{(\alpha + i\beta)x} = e^{\alpha x} \cos \beta x + ie^{\alpha x} \sin \beta x\) , so, by (28)

\[
\begin{array}{l} \mathrm{D} \left(\mathcal {R} \left(\frac {1}{\alpha + i \beta} e ^ {(\alpha + i \beta) x}\right)\right) = e ^ {\alpha x} \cos \beta x \\ \mathrm{D} \left(\mathcal {I} \left(\frac {1}{\alpha + i \beta} e ^ {(\alpha + i \beta) x}\right)\right) = e ^ {\alpha x} \sin \beta x. \end{array}
\]

In the same way one reduces the evaluation of a primitive of \(x^n e^{\alpha x} \cos \beta x\) , or of \(x^n e^{\alpha x} \sin \beta x\) ( \(n\) an integer \(> 0\) ) to that of a primitive of \(x^n e^{(\alpha + i\beta)x}\) ; now, the formula for integration by parts of order \(n + 1\) (II, p.~60, formula (11)) shows that a primitive of this last function is

\[
e ^ {(\alpha + i \beta) x} \left[ \frac {x ^ {n}}{\alpha + i \beta} - \frac {n x ^ {n - 1}}{(\alpha + i \beta) ^ {2}} + \frac {n (n - 1) x ^ {n - 2}}{(\alpha + i \beta) ^ {3}} + \dots + (- 1) ^ {n} \frac {n !}{(\alpha + i \beta) ^ {n + 1}} \right].
\]

By Euler's formulae one can on the other hand express every positive integral power of \(\cos x\) or of \(\sin x\) as a linear combination of exponentials \(e^{ipx}\) ( \(p\) a positive or negative integer). By formula (28) one can thus express a primitive of a function of the form \(x^n e^{\alpha x}(\cos \beta x)^r (\sin \gamma x)^s\) as a linear combination of functions of the form \(x^p e^{\alpha x}\cos \lambda x\) and \(x^p e^{\alpha x}\sin \mu x\) and \((n, p, r, s\) integers, \(\alpha, \beta, \gamma, \lambda, \mu\) real).

Example. One has

\[
\sin^ {2 n} x = \frac {(- 1) ^ {n}}{2 ^ {2 n}} \left(e ^ {i x} - e ^ {- i x}\right) ^ {2 n} = \frac {(- 1) ^ {n}}{2 ^ {2 n}} \left(e ^ {2 n i x} - \binom {2 n} {1} e ^ {(2 n - 2) i x} + \dots + e ^ {- 2 n i x}\right)
\]

whence

\[
\begin{array}{l} \int_ {0} ^ {x} \sin^ {2 n} t d t = \frac {(- 1) ^ {n}}{2 ^ {2 n}} \left(\frac {1}{n} \sin 2 n x - \binom {2 n} {1} \frac {1}{n - 1} \sin (2 n - 2) x + \dots \right. \\ \qquad \qquad \qquad \qquad \qquad \qquad \qquad \qquad \qquad \qquad + (- 1) ^ {n - 1} \binom {2 n} {n - 1} \sin 2 x + (- 1) ^ {n} \binom {2 n} {n} x \Bigg) \end{array}
\]

and in particular

\[
\int_ {0} ^ {\pi / 2} \sin^ {2 n} t   d t = \binom {2 n} {n} \frac {1}{2 ^ {2 n}}   \frac {\pi}{2} = \frac {1 . 3 . 5 \ldots (2 n - 1)}{2 . 4 . 6 \ldots 2 n}   \frac {\pi}{2}.\tag{29}
\]

